\documentclass[12pt]{article}
\usepackage[utf8]{inputenc}
\usepackage[T1]{fontenc}
\usepackage{hyperref}
\title{Semandang Language Resources for K-12 Curriculum Development}
\author{Compiled by Claude Sonnet 5.5 (claude-sonnet-5-5) for Andr\'as Kornai}
\date{}
\begin{document}
\maketitle

\section{Introduction}
Semandang (also called Onya Darat; ISO 639-3 sdq) is a Southern Land Dayak Austronesian language of Ketapang regency, West Kalimantan, Indonesia. Sources estimate roughly 10,000 to 50,000 speakers. Searches in English and Indonesian found no textbook, dictionary or school material, so this dossier is short.

\section{Primary Educational Resources}
No primers, textbooks or local-content (muatan lokal) materials for Semandang were found. The reference description says it is taught as a subject in some schools, but no source confirming this turned up.

\section{Linguistic and Reference Materials}
\begin{itemize}
\item Wikipedia article on Semandang. \url{https://en.wikipedia.org/wiki/Semandang_language}
\item Wikipedia article on the Land Dayak languages. \url{https://en.wikipedia.org/wiki/Land_Dayak_languages}
\item MPI-EVA project page, ``Documentation of the phonetic structures of Onya Darat''. \url{https://www.eva.mpg.de/linguistics/past-research-resources/documentation-and-description/documentation-of-the-phonetic-structures-of-onya-darat/}
\item Project member page for the same MPI-EVA project. \url{https://www.eva.mpg.de/linguistics/past-research-resources/phonetics-and-phonology/documentation-of-the-phonetic-structures-of-onya-darat}
\item Glottolog entry for ISO code sdq. \url{https://glottolog.org/resource/languoid/iso/sdq}
\item Survey document on Land Dayak languages of West Kalimantan (Soriente and Inagaki). \url{https://lingdy.aa-ken.jp/wp-content/uploads/2012/01/120217_soriente_inagaki_h2.pdf}
\end{itemize}
A secondary source says the only academic publication on Onya Darat is a short article on its orthography (Tadmor 2009). No URL for it was found.

\section{Religious and Community Resources}
\begin{itemize}
\item Joshua Project page on Semandang. \url{https://joshuaproject.net/people_groups/14838}
\item Global Recordings Network page listed as Semandang. It uses a different code (sdm), so it may describe another language and should be checked before use. \url{https://globalrecordings.net/en/language/sdm}
\end{itemize}

\section{Digital and Audio Resources}
No Semandang-specific corpus, audio or software was confirmed.

\section{Curriculum Development Context}
Joshua Project notes that younger speakers are shifting to Indonesian and Malay. Indonesian local-content curricula allow regional languages in schools, but nothing Semandang-specific was found. A curriculum would need new materials built from the MPI-EVA documentation and from community speakers.
\end{document}